\chapter{اللقاء التاسع والثلاثون: وعد الله ووعد الشيطان والحكمة وبداية آيات الربا}

\section{مقدمة}
السلام عليكم ورحمة الله وبركاته. الحمد لله رب العالمين، وأشهد أن لا إله إلا الله وحده لا شريك له، وأشهد أن محمداً عبده ورسوله، صلى الله عليه وعلى آله وصحبه وسلم.

استأنف الشيخ هذا اللقاء من قوله تعالى: \begin{ayah}الشَّيْطَانُ يَعِدُكُمُ الْفَقْرَ وَيَأْمُرُكُمْ بِالْفَحْشَاءِ وَاللَّهُ يَعِدُكُمْ مَغْفِرَةً مِنْهُ وَفَضْلًا\end{ayah}، ثم انتقل إلى آيات الحكمة والتذكر والنفقة والصدقات، حتى دخل في أوائل آيات الربا وما فيها من الوعيد الشديد والزجر البالغ.

وفي هذا المجلس تظهر عناية \rawi{الطبري} بجمع الأقوال المتقاربة تحت أصل جامع، وربط خواتيم الآيات بأسمائه سبحانه، وبيان أن التشريع كله قائم على العلم والحكمة والعدل، لا على محض التفريق بين المتماثلات كما زعمه أهل الباطل.

\separator

\section{وعد الشيطان ووعد الله}
\isnad{قال أبو جعفر:} يعني تعالى ذكره بقوله: \begin{ayah}الشَّيْطَانُ يَعِدُكُمُ الْفَقْرَ\end{ayah} أي يخوفكم بالفاقة إذا أنفقتم، ويأمركم بالمعاصي وترك ما أوجب الله عليكم. وأما قوله: \begin{ayah}وَاللَّهُ يَعِدُكُمْ مَغْفِرَةً مِنْهُ وَفَضْلًا\end{ayah}، فمعناه أنه يعد عباده ستر الذنوب ومحو آثارها، ويعدهم كذلك بالخلف والبركة والرزق.

ووقف الشيخ عند هذا التقابل وقفة تربوية جليلة: فالشيطان لا يواجه العبد عادة بصورة مجردة من الشر، بل يلبس له المعصية ثوب المصلحة، ويصور له الطاعة كأنها خصم لمنافعه أو سبب في ضياعها. فمن همّ بالصدقة خوفه الفقر، ومن همّ بترك المعصية منّاه بما يفوته من لذة أو ربح أو جاه.

وقد قرأ الشيخ في هذا الموضع حديث \begin{hadith}إن للشيطان لمة من ابن آدم، وللملك لمة\end{hadith}، وبيّن أن الخطر ليس في مجرد ورود الوسوسة، بل في الاستجابة لها والانقياد بعدها. فالعبد قد تعرض له الخواطر، لكن الشر الحقيقي أن يتبعها، ولهذا كان الاستعاذة بالله طلباً للعصمة من الاستجابة، لا من أصل الخاطر الذي لا يملك دفعه دائماً.

\begin{fawaid}[الشيطان يقلب الموازين في نظر العبد]
من أعظم مداخل الشيطان أن يجعل الطاعة كأنها مفسدة عاجلة، والمعصية كأنها منفعة قريبة، فيحتاج المؤمن دائماً إلى تجديد اليقين بأن وعد الله حق، وأن ما يأمر به الشيطان غرور وخداع.
\end{fawaid}

\section{أسماء الله في خواتيم الآيات من مفاتيح الفهم}
عند قوله تعالى: \begin{ayah}وَاللَّهُ وَاسِعٌ عَلِيمٌ\end{ayah} نبّه الشيخ إلى أصل متكرر في هذه الدروس: أن ختام الآية باسم من أسماء الله ليس فاصلة لفظية مجردة، بل هو جزء من الهداية والمعنى.

فاسم \begin{ayah}وَاسِعٌ\end{ayah} هنا يورث المؤمن الثقة بسعة فضل الله وخزائنه، فلا يتوهم أن النفقة تنقصه نقصاً لا يعوض. واسم \begin{ayah}عَلِيمٌ\end{ayah} يفتح باب المراقبة والرجاء معاً؛ فهو سبحانه يعلم ما يخرجه العبد، ويعلم قصده فيه، ويحصيه له، ثم يجازيه عليه.

ومن هنا أكد الشيخ أن باب العلم بأسماء الله ومحامده يدخل في جميع أبواب الإسلام: في فهم الأحكام، وفي الإيمان بالقدر، وفي الثقة بوعد الله، وفي حسن تلقي أوامره ونواهيه.

\begin{fawaid}[خواتيم الآيات تربي القلب على تلقي الحكم]
كلما ازداد العبد عناية بالأسماء الإلهية التي تختم الآيات ازداد فهماً للحكم الذي قبلها، لأن الاسم المختار في الخاتمة يدل على المعنى المراد تقريره في القلب والعمل.
\end{fawaid}

\section{معنى الحكمة عند الطبري}
في قوله تعالى: \begin{ayah}يُؤْتِي الْحِكْمَةَ مَنْ يَشَاءُ\end{ayah} نقل \rawi{الطبري} أقوال السلف: فقيل هي القرآن والفقه به، وقيل العلم بالدين، وقيل الفهم، وقيل الخشية، وقيل النبوة. ثم جمع هذه الأقوال كلها تحت أصل واحد، فقال: الحكمة هي الإصابة في القول والفعل.

وبيّن الشيخ جمال هذا المسلك؛ فإن \rawi{الطبري} لا يعامل اختلاف السلف هنا على أنه تضاد، بل يرد الجميع إلى المعنى الجامع، ثم يجعل كل قول منها مثالاً لصنف من أصناف الحكمة. فالنبوة أعلاها، والعلم بالكتاب من أجلها، والخشية أثر من آثارها، والفقه في الدين من لوازمها.

وتوقف الشيخ طويلاً مع كلمة \rawi{مالك}: «المعرفة بالدين والفقه فيه والاتباع له»، وعدّها من أنفس ما قيل في الباب؛ لأن العلم إذا لم يثمر اتباعاً لم يكتمل وصف الحكمة في صاحبه.

\begin{fawaid}[الحكمة معنى جامع لا كلمة ضيقة]
من دقائق قراءة الطبري أنه يجمع الألفاظ العامة في القرآن إلى أصولها الكلية، ثم يبين أن الأقوال المأثورة فيها أمثلة وصور، لا خصومات متباينة يلزم إسقاط بعضها ببعض.
\end{fawaid}

\section{التذكر ليس مجرد العلم}
فسر \rawi{الطبري} قوله تعالى: \begin{ayah}وَمَا يَذَّكَّرُ إِلَّا أُولُو الْأَلْبَابِ\end{ayah} بأنهم الذين يتعظون بوعد الله ووعيده، فينزجرون عما نهوا عنه، ويطيعون فيما أمروا به.

وهنا نبه الشيخ إلى فرق مهم: فالقرآن هدى للناس من جهة البيان والبذل، لكن الذي ينتفع به حقيقة هم الذين يقبلونه ويعملون به. فالتذكر في مثل هذه الآية ليس مجرد الاستحضار الذهني أو المعرفة المجردة، بل هو تذكر يثمر امتثالاً واتعاظاً وانقياداً.

ولهذا كان من أوصاف أولي الألباب أنهم إذا عرض عليهم الحق قبلوه وساروا معه، لا أنهم عرفوه ثم وقفوا عند حد الإدراك النظري.

\begin{fawaid}[علامة العقل النافع قبول الحق بعد ظهوره]
ليس العاقل في ميزان القرآن من عرف المعنى فقط، بل من حمله علمه على الانقياد، ولهذا خص الله أولي الألباب بحسن التذكر والاتعاظ.
\end{fawaid}

\section{علم الله بالنفقة والنذر}
عند قوله تعالى: \begin{ayah}وَمَا أَنْفَقْتُمْ مِنْ نَفَقَةٍ أَوْ نَذَرْتُمْ مِنْ نَذْرٍ فَإِنَّ اللَّهَ يَعْلَمُهُ\end{ayah}، بيّن \rawi{الطبري} أن الله لا يخفى عليه شيء من نفقات العباد ونذورهم، بل يحصيها ويجازي عليها بحسب النية والقصد.

وذكر الشيخ أن العلم هنا يتضمن معناه الأصلي، ويتضمن أيضاً لازمَه من الجزاء والحساب؛ فالله يعلم العمل، ويعلم ما قارنه من الإخلاص أو الرياء، ثم يرتب على ذلك ثوابه أو عقابه. فليست النفقة سواءً إذا اتحد ظاهرها واختلف قصدها.

ومن هنا جاء التهديد في آخر الآية: \begin{ayah}وَمَا لِلظَّالِمِينَ مِنْ أَنْصَارٍ\end{ayah}، أي الذين وضعوا النفقة أو النذر في غير موضعهما: رياءً للناس، أو طاعة للشيطان، أو صرفاً للمال في غير ما يحبه الله.

\section{إبداء الصدقات وإخفاؤها}
في قوله تعالى: \begin{ayah}إِنْ تُبْدُوا الصَّدَقَاتِ فَنِعِمَّا هِيَ وَإِنْ تُخْفُوهَا وَتُؤْتُوهَا الْفُقَرَاءَ فَهُوَ خَيْرٌ لَكُمْ\end{ayah}، رجح \rawi{الطبري} أن الآية عامة في الصدقات، إلا أن الزكاة الواجبة يبقى إظهارها في الجملة أليق بسائر الفرائض، أما صدقة التطوع فإخفاؤها أفضل وأبعد عن الرياء.

وأشار الشيخ إلى أن هذا من المواضع التي يجمع فيها \rawi{الطبري} بين ظاهر العموم وبين ما انعقد عليه المعنى الشرعي في تمييز الواجب من التطوع. كما مر في أثناء الدرس ذكر اختلاف القراءات في قوله: \begin{ayah}وَيُكَفِّرُ عَنْكُمْ مِنْ سَيِّئَاتِكُمْ\end{ayah}، مع بقاء أصل المعنى متقارباً: أن الصدقة من أسباب تكفير الذنوب.

\begin{fawaid}[إخفاء التطوع أقرب إلى الإخلاص وإظهار الفرض أقرب إلى البيان]
ليست جميع الصدقات في مرتبة واحدة من جهة الإعلان والإسرار؛ فباب التطوع أليق به الخفاء غالباً، أما الواجب فيغلب فيه معنى الإظهار والامتثال الظاهر وإقامة الشعيرة.
\end{fawaid}

\separator

\section{بداية آيات الربا: القيام يوم القيامة ودعوى المماثلة}
دخل الشيخ بعد ذلك في قوله تعالى: \begin{ayah}الَّذِينَ يَأْكُلُونَ الرِّبَا لَا يَقُومُونَ إِلَّا كَمَا يَقُومُ الَّذِي يَتَخَبَّطُهُ الشَّيْطَانُ مِنَ الْمَسِّ\end{ayah}، وذكر تفسير \rawi{الطبري} لهذا القيام بأنه حالهم عند البعث من قبورهم، يبعثون على صورة فظيعة تناسب جرمهم وعظيم جنايتهم.

ثم نبه \rawi{الطبري} إلى معنى دقيق: أن التعبير بالأكل لا يراد به قصر التحريم على من أكل الربا في مطعمه، بل هو وصف لأشهر وجوه انتفاعهم به، وإلا فالوعيد يتناول العامل به، والآخذ له، والمعطي، وكل من دخل في سببه عالماً به.

وأصل جريمتهم أن قالوا: \begin{ayah}إِنَّمَا الْبَيْعُ مِثْلُ الرِّبَا\end{ayah}، فزعموا أن الزيادة في البيع كالزيادة الربوية عند حلول الأجل. فرد الله عليهم بقوله: \begin{ayah}وَأَحَلَّ اللَّهُ الْبَيْعَ وَحَرَّمَ الرِّبَا\end{ayah}.

\begin{fawaid}[التعبير بالأكل في الربا خرج على الأغلب لا على الحصر]
من دخل في الربا كسباً أو كتابة أو إعانة أو أخذاً أو إعطاءً لا يخرج من الوعيد بحجة أنه لم يأكل المال بعينه، لأن الآية خرجت على أشهر صور الانتفاع لا على قصر الحكم عليها.
\end{fawaid}

\section{التسليم لحكم الله لا ينافي طلب حكمته}
وقف الشيخ هنا عند مسألة مهمة جداً: هل قوله تعالى \begin{ayah}وَأَحَلَّ اللَّهُ الْبَيْعَ وَحَرَّمَ الرِّبَا\end{ayah} يفيد أن الأحكام مجرد مشيئة لا حكمة وراءها؟ وقرر بوضوح أن هذا فهم باطل.

فالآية رد على من زعم التماثل بين البيع والربا عند الله، لا على أصل أن أحكام الله كلها جارية على العلم والعدل والرحمة. فنحن نتلقى التشريع أولاً بالتسليم؛ لأن الله لا معقب لحكمه، ثم نزداد يقيناً بأن كل ما أحله أو حرمه فبحكمة، علمناها أو جهلناها.

وبيّن الشيخ أن من عرف حقيقة الربا وآثاره في الأفراد والمجتمعات ازداد فهماً لشيء من حكمة التحريم، لكن بقاء بعض الحكم خافياً على بعض الناس لا يقدح في أصل العلم بأن شريعة الله قائمة على الكمال.

\begin{fawaid}[إثبات الحكمة في التشريع لا يعارض كمال التسليم]
المؤمن يسلم لحكم الله لأنه حكم الله، ويوقن في الوقت نفسه أن حكمه كله عدل ورحمة وعلم، فلا يجعل التسليم ذريعة لنفي الحكمة، ولا طلب الحكمة ذريعة لتعليق الطاعة عليها.
\end{fawaid}

\section{فمن جاءه موعظة من ربه فانتهى}
فسر \rawi{الطبري} الموعظة هنا بالتذكير والتخويف والتحريم الذي جاء في القرآن، فمن بلغه هذا البيان فانتهى عن الربا فله ما سلف قبل التحريم، وأمره بعد ذلك إلى الله في تثبيته وعصمته.

وأفاض الشيخ في تنزيل هذا المعنى على واقع الناس والأسواق، فذكر أن كثيراً من المعاملات قد يدخلها الربا من حيث لا يشعر أصحابها، إما لجهلهم بالحكم، وإما لجهلهم بحقيقة الصورة التي يتعاملون بها. ومن هنا كان الواجب على التجار أن يسألوا قبل الإقدام، وعلى أهل العلم أن يواكبوا صور المعاملات المستجدة ولا يتركوا الناس يتخبطون في الأسواق بلا بيان.

فإذا قامت الموعظة وظهر الحكم، لم يبق لأحد أن يقول: لم أكن أعلم، وهو قادر على السؤال والبحث وطلب البيان.

\begin{fawaid}[من جاءته الموعظة لزمه الوقوف عندها]
العذر بالجهل لا يفتح باب التهاون في السؤال، بل يوجب على المسلم أن يتعلم ما يدخل فيه من أبواب المعاملات، فإذا بان له الحق لم يجز له أن يعود إلى ما نهاه الله عنه.
\end{fawaid}

\section{يمحق الله الربا ويربي الصدقات}
في قوله تعالى: \begin{ayah}يَمْحَقُ اللَّهُ الرِّبَا وَيُرْبِي الصَّدَقَاتِ\end{ayah}، بيّن \rawi{الطبري} أن المحق هو النقص والإذهاب، وأن الله وإن وسع لبعض أهل الربا في الظاهر زمناً، فإن عاقبة هذا المال إلى ذلة أو تلف أو ذهاب بركته. وأما الصدقة فإن الله ينميها لصاحبها ويضاعف أجرها في الدنيا والآخرة.

وقرأ الشيخ حديث تربية الله للصدقة لصاحبها حتى تصير كالشيء العظيم، وجعل هذا من أجمل أمثلة تفسير القرآن بالسنة؛ إذ يبين الحديث ما أجملته الآية، ويقرب المعنى إلى القلب.

ثم جاء الختام: \begin{ayah}وَاللَّهُ لَا يُحِبُّ كُلَّ كَفَّارٍ أَثِيمٍ\end{ayah}، فبيّن \rawi{الطبري} أن هذا الوصف متجه إلى من أصر على استحلال الربا والمقام عليه، لا إلى من وقع فيه ثم رجع وتاب.

\section{ذروا ما بقي من الربا فإن لم تفعلوا فأذنوا بحرب}
بلغ الوعيد غايته في قوله تعالى: \begin{ayah}يَا أَيُّهَا الَّذِينَ آمَنُوا اتَّقُوا اللَّهَ وَذَرُوا مَا بَقِيَ مِنَ الرِّبَا\end{ayah}، فبيّن \rawi{الطبري} أن المراد ترك ما بقي من الزيادات الربوية التي لم تقبض بعد، والاكتفاء برؤوس الأموال دون ما زاد عليها.

ثم جاء قوله تعالى: \begin{ayah}فَإِنْ لَمْ تَفْعَلُوا فَأْذَنُوا بِحَرْبٍ مِنَ اللَّهِ وَرَسُولِهِ\end{ayah}، فذكر \rawi{الطبري} أن القراءة الأقوى عنده بمعنى: اعلموا واستيقنوا بحرب من الله ورسوله. وهذا من أشد ما ورد في الوعيد في أبواب المعاملات، ولذلك شدد الشيخ في التحذير من التهاون في القروض الربوية والصور المعاصرة التي توسع فيها كثير من الناس والفتاوى.

ثم ختم المجلس بقوله تعالى: \begin{ayah}وَإِنْ تُبْتُمْ فَلَكُمْ رُءُوسُ أَمْوَالِكُمْ\end{ayah}، أي إن تبتم من الربا فلكم أصل المال فقط دون الزيادة المحرمة.

\begin{fawaid}[باب الربا من أشد أبواب الوعيد في الشريعة]
إذا كان الله قد توعد آكلي الربا والمقيمين عليه بالحرب، علم المؤمن أن هذا الباب ليس من صغائر المخالفات ولا من المسائل التي يتوسع فيها طلباً للدنيا، بل هو من أعظم ما يجب الاحتياط فيه للدين والذمة.
\end{fawaid}

\separator

خلاصة هذا اللقاء أن السورة نقلت المؤمن هنا من الحديث عن الإنفاق والصدقة ووعد الله بالخلف إلى أصلين عظيمين: أصل الحكمة التي يهبها الله لمن يشاء، وأصل التمييز بين الكسب الطيب الذي يباركه الله والكسب الخبيث الذي يمحقه ولو كثر. وفي هذا المقطع يظهر توازن بديع بين الترغيب والترهيب: وعد الله بالمغفرة والفضل، ثم وعيده الشديد في الربا لمن عاند وأصر.
